\subsection*{定理 5.7.}
\underline{在 5.6 的假设下, 置}
\[
\begin{aligned}
E^{n-2}(\Delta,\mathbb Q_\ell(j))
&=\underline{\text{在 }}
H^{n-2}(\Delta,\mathbb Q_\ell(j))\\
&\qquad\underline{\text{中与下列空间之像正交的部分: }}
H^{n-2}(X,\mathbb Q_\ell(j)).
\end{aligned}
\]
\underline{(这是 $\Delta$ 的上同调的\textquotedblleft 消灭\textquotedblright{}部分; 参见 5.4.4.) 则有直和分解} (参见 4.2)
\begin{equation}\tag{5.7.1}
E_2^{1,n-1}
=H^1(\mathbb P^1,R^{n-1}\rho_*\mathbb Q_\ell)
\xleftarrow[\;f^*+i_*g^*\;]{\sim}
\operatorname{Prim}^n(X,\mathbb Q_\ell)
\oplus E^{n-2}(\Delta,\mathbb Q_\ell(-1)).
\end{equation}

\underline{证明.}
由退化性,
\[
E_2^{1,n-1}\simeq
F^1H^n(\widetilde X,\mathbb Q_\ell)
\big/F^2H^n(\widetilde X,\mathbb Q_\ell).
\]
由 (A),
\begin{equation}\tag{5.7.2}
F^1H^n(\widetilde X,\mathbb Q_\ell)
=\ker\!\left(
k^*:H^n(\widetilde X,\mathbb Q_\ell)\longrightarrow
H^n(X_{\bar\eta},\mathbb Q_\ell)
\right),
\end{equation}
由 (5.1.2), 它对应于
\[
\ker\!\left(
H^n(X,\mathbb Q_\ell)
\oplus H^{n-2}(\Delta,\mathbb Q_\ell(-1))
\xrightarrow{\;m^*+h_*\;}
H^n(X_{\bar\eta},\mathbb Q_\ell)
\right).
\]
在 (LV) 下, $k^*$ 也可写为
\begin{equation}\tag{5.7.3}
\operatorname{Prim}^n(X,\mathbb Q_\ell)
\oplus H^{n-2}(X,\mathbb Q_\ell(-1))
\oplus H^{n-2}(\Delta,\mathbb Q_\ell(-1))
\xrightarrow{\;m^*+m^*L+h_*\;}
H^n(X_{\bar\eta},\mathbb Q_\ell),
\end{equation}
且中间的箭头
\begin{equation}\tag{5.7.4}
H^{n-2}(X,\mathbb Q_\ell(-1))
\xrightarrow{\;m^*L\;}
H^n(X_{\bar\eta},\mathbb Q_\ell)
\end{equation}
为同构 (参见 5.2.3 和 (5.2.2.3)). 因而, 以下式表示到本原部分的投影:
\begin{equation}\tag{5.7.5}
\wp:H^n(X,\mathbb Q_\ell)\longrightarrow
\operatorname{Prim}^n(X,\mathbb Q_\ell)
\end{equation}
我们得到
\footnote{印本在 (5.7.6) 的标记中在 $g_*$ 后写作 $2^*$. 此处直接使用爆破分解的逆 (4.2.2), 因而要求 $i^*$.}
\begin{equation}\tag{5.7.6}
F^1H^n(\widetilde X,\mathbb Q_\ell)
\xrightarrow[\;\wp f_*\oplus(-g_*i^*)\;]{\sim}
\operatorname{Prim}^n(X,\mathbb Q_\ell)
\oplus H^{n-2}(\Delta,\mathbb Q_\ell(-1)).
\end{equation}
由 (5.6.2.2) 及退化性,
\begin{equation}\tag{5.7.7}
F^2H^n(\widetilde X,\mathbb Q_\ell)
\simeq E_2^{2,n-2}
\simeq H^{n-2}(X_{\bar\eta},\mathbb Q_\ell(-1)).
\end{equation}

另一方面, $F^2H^n(\widetilde X,\mathbb Q_\ell)=H^{n-2}(X_{\bar\eta},\mathbb Q_\ell(-1))$ 到 $H^n(\widetilde X,\mathbb Q_\ell)$ 的包含映射为 Gysin 映射 $k_*$.
\footnote{印本在首句中写作 $F^2H^N$, 并略去了所列 Gysin 映射源中的 Tate 扭转 $(-1)$. 周围的对应以及 Gysin 映射使次数移动二这一事实, 都要求此处所给的写法.}
严格说来, Gysin 映射仅定义为映射
\[
H^{i-2}(X_{\bar\eta},\mathbb Q_\ell(-1))
\longrightarrow
H^i(\widetilde X\times_k k(\bar\eta),\mathbb Q_\ell),
\]
但由于典范同构 (SGA 4 XII 5.4)
\[
H^i(X,\mathbb Q_\ell)\xrightarrow{\sim}
H^i(X\otimes_kK,\mathbb Q_\ell)
\]
对每个代数闭域 $K/k$ 都成立, 希望读者允许我们这样的滥用记号.
\begin{editorialnote}
原印本及英译底本在上列严格意义的 Gysin 映射靶中写作 $X\times_k k(\bar\eta)$. 这里 $k$ 是一般纤维到 $\widetilde X$ 的浸入, 其 Gysin 映射的靶应为 $\widetilde X\times_k k(\bar\eta)$ 的上同调; 此处补回波浪号. 后一条代数闭域扩张下的上同调不变性则适用于任意固有概形, 特别适用于 $\widetilde X$.
\end{editorialnote}

由 (5.1.3), Gysin 映射表示为
\begin{equation}\tag{5.7.8}
H^{n-2}(X_{\bar\eta},\mathbb Q_\ell(-1))
\xrightarrow{\;m_*\oplus(-h^*)\;}
H^n(X,\mathbb Q_\ell)
\oplus H^{n-2}(\Delta,\mathbb Q_\ell(-1)).
\end{equation}
由弱 Lefschetz 定理,
\[
m_*H^{n-2}(X_{\bar\eta},\mathbb Q_\ell(-1))
=m_*m^*H^{n-2}(X,\mathbb Q_\ell(-1))
=L_XH^{n-2}(X,\mathbb Q_\ell(-1))
\]
为本原部分的正交补. 因而有交换图
\begin{equation}\tag{5.7.9}
\begin{tikzcd}[column sep=huge,row sep=huge]
F^2H^n(\widetilde X,\mathbb Q_\ell)
  \arrow[r,hook]
& F^1H^n(\widetilde X,\mathbb Q_\ell)
  \arrow[d,"\wp f_*\oplus(-g_*i^*)","\wr"']
\\
H^{n-2}(X_{\bar\eta},\mathbb Q_\ell(-1))
  \arrow[u,"k_*"',"\wr"]
  \arrow[r,"0\oplus(-h^*)"]
& \operatorname{Prim}^n(X,\mathbb Q_\ell)
  \oplus H^{n-2}(\Delta,\mathbb Q_\ell(-1)).
\end{tikzcd}
\end{equation}
\begin{editorialnote}
原印本及英译底本在图 (5.7.9) 底边写作 $0\oplus h^*$. 由 (5.7.8), $k_*$ 在爆破分解下为 $m_*\oplus(-h^*)$, 而 $\wp m_*=0$, 所以交换性要求底边为 $0\oplus(-h^*)$. 此处订正该符号.
\end{editorialnote}

这就证明了所需的 (5.7.1).

\subsection*{注 5.7.10.}
在 5.7 中, 假设 (A) 可以换成较弱的假设: 对 $q=n$, 映射
\begin{equation}\tag{5.7.11}
H^q(\widetilde X,\mathbb Q_\ell)\longrightarrow
H^0(\mathbb P^1,R^q\rho_*\mathbb Q_\ell)
\end{equation}
以及
\begin{equation}\tag{5.7.12}
H^q(\widetilde X,\mathbb Q_\ell)\longrightarrow
H^q(X_{\bar\eta},\mathbb Q_\ell)
\end{equation}
具有相同的核, 正如我们的证明所表明的. 若无此条件, 结论可能不成立; 参见 5.9.
